% Solutions manual front matter: solutions PDF pp. 1--2.
\noindent\emph{Solutions to Axler's Linear Algebra Done Right, 4th Edition}, by
blanketism. Began July 4, 2026; last edited July 20, 2026.

\section*{Preface}\addcontentsline{toc}{section}{Preface}

\emph{Note: For any corrections, suggestions, or feedback regarding this
solutions manual, readers are encouraged to reach out through the following
channels:}
\begin{itemize}
\item Discord Errata Server: \url{https://discord.gg/qecQk37vfm}
\item Reddit: \url{https://www.reddit.com/user/ansv9a8fdh3/}
\item GitHub: \url{https://github.com/blanketism/Axler-Solutions}
\end{itemize}

This solutions manual accompanies the fourth edition of Sheldon Axler's
\emph{Linear Algebra Done Right}. While several community-maintained solution
sets exist for the third edition, the fourth edition introduces enough new and
restructured material to warrant a dedicated resource--one written from scratch
with the new text in mind.

This manual is written primarily for undergraduates encountering rigorous
linear algebra for the first time. Axler's text assumes little beyond
mathematical maturity and a basic familiarity with proof writing. There are
some solutions in sections that may have mathematical identities that require
additional background or insight, and these are explained as needed to ensure
clarity and completeness.

A reader may notice certain notes underneath each exercise statement. A guide
is as follows:
\begin{itemize}
\item \textbf{Axler Hint.} \emph{These are hints (explicitly indicated, or
  otherwise) provided by Axler to guide the reader and assist in solving the
  exercise.}
\item \textbf{Axler Note.} \emph{Small excerpts written by Axler that either
  provide a reference to some term used in the problem, or offer additional
  context or clarification relevant to understanding the exercise. I have
  abstained from including every note provided.}
\item \textbf{Manual Remark.} \emph{These are remarks provided by myself to
  assist the reader in understanding the exercise, offering additional
  insights, clarifications, or alternative perspectives that may not be
  immediately apparent from the problem statement alone.}
\end{itemize}

The goal here is not to replace the work of reading and thinking. Linear
algebra, particularly in Axler's determinant-free, operator-theoretic style,
rewards patience and independent struggle in ways that are difficult to
replicate by reading someone else's solution. This manual is best used after a
genuine attempt at each problem, as a means of checking reasoning, identifying
gaps, or becoming unstuck--not as a first resort.

Despite best efforts, mistakes are inevitable in any document of this scope:
incorrect arguments, incomplete cases, notational slips, and the occasional
solution that is simply wrong. It has been some time since I first worked
through Axler's text, and while I have made every effort to ensure accuracy, I
recognize that errors may still exist. Readers who encounter such mistakes are
encouraged to report them through the channels listed above, and all
contributions to improving this manual are greatly appreciated.

Finally, I am grateful to Professor Axler for writing a transformative text. I
recall my first experience with this text as an undergraduate, and it
profoundly shaped my understanding and appreciation of linear algebra. This
solutions manual is my way of giving back and supporting future students in
their journey through Axler's material. I hope that it serves as a helpful
companion, providing clarity and guidance while encouraging readers to engage
deeply with the subject on their own terms.

\subsection*{In this edition}

The manual covers Chapters 1--4 and Sections 5A--5B. Its solutions are printed
below, grouped by section; in the text, each solved exercise is followed by its
solution. The alternate proofs from the manual's appendix are placed at the end
of the corresponding solutions. The manual's restatements of the exercises are
omitted, since the exercises appear in the text.
